
\documentclass{article}
\usepackage{colortbl}
\usepackage{makecell}
\usepackage{multirow}
\usepackage{supertabular}

\begin{document}

\newcounter{utterance}

\centering \large Interaction Transcript for game `cladder', experiment `full\_v1.5\_default', episode 4119 with reflexion{-}Qwen3.5{-}9B.
\vspace{24pt}

{ \footnotesize  \setcounter{utterance}{1}
\setlength{\tabcolsep}{0pt}
\begin{supertabular}{c@{$\;$}|p{.15\linewidth}@{}p{.15\linewidth}p{.15\linewidth}p{.15\linewidth}p{.15\linewidth}p{.15\linewidth}}
    \# & \multicolumn{2}{c}{Player} && \multicolumn{2}{c}{Game Master} \\
    \hline

    \theutterance \stepcounter{utterance}  
    & & & \multicolumn{4}{p{0.6\linewidth}}{
        \cellcolor[rgb]{0.9,0.9,0.9}{
            \makecell[{{p{\linewidth}}}]{
                \texttt{\tiny{[P1$\langle$GM]}}
                \texttt{You have attempted this same task before and did not succeed. These are the plans you wrote after those attempts.} \\
\texttt{Trial \#1: New plan: I incorrectly concluded "No" despite my own reasoning explicitly stating that the probability drops from 47\% to 15\%, which mathematically confirms a decrease. The error was a failure to align my final conclusion with the evidence I generated. Next time, I will ensure that after calculating the probabilities and identifying the direction of the effect, I explicitly verify that my opening "Yes" or "No" matches the calculated outcome before finalizing the response. Specifically, since \$P(\textbackslash{}text\{Curly\} | \textbackslash{}text\{Citrus\}) < P(\textbackslash{}text\{Curly\} | \textbackslash{}neg \textbackslash{}text\{Citrus\})\$, the answer must be "Yes".} \\
\\ 
\\ 
\texttt{You are an expert in causal inference. The following question is not a typical commonsense query, but rather a meticulously designed question created by a professor specializing in causal inference, intended to assess the students' mastery of the course content.} \\
\texttt{Imagine a self{-}contained, hypothetical world with only the following conditions, and without any unmentioned factors or causal relationships: Eating citrus has a direct effect on vitmain C. Vitmain C has a direct effect on curly hair. For patients not consuming citrus, the probability of curly hair is 47\%. For patients consuming citrus, the probability of curly hair is 15\%. Will citrus intake decrease the chance of curly hair?} \\
\texttt{Start your answer with 'Yes' or 'No', followed by additional reasoning or evidence to support your explanation.} \\
            }
        }
    }
     \\ \\

    \theutterance \stepcounter{utterance}  
    & \multicolumn{4}{p{0.6\linewidth}}{
        \cellcolor[rgb]{0.561,0.737,0.561}{
            \makecell[{{p{\linewidth}}}]{
                \texttt{\tiny{[P1$\rangle$GM]}}
                \texttt{Yes, citrus intake decreases the chance of curly hair.} \\
\\ 
\texttt{In this hypothetical world, we are given the conditional probabilities directly:} \\
\texttt{{-} \$P(\textbackslash{}text\{Curly\} | \textbackslash{}neg \textbackslash{}text\{Citrus\}) = 47\textbackslash{}\%\$} \\
\texttt{{-} \$P(\textbackslash{}text\{Curly\} | \textbackslash{}text\{Citrus\}) = 15\textbackslash{}\%\$} \\
\\ 
\texttt{The causal chain is defined as: Citrus \$\textbackslash{}rightarrow\$ Vitamin C \$\textbackslash{}rightarrow\$ Curly Hair. Since there are no unmentioned factors or confounders in this self{-}contained world, the observed conditional probabilities reflect the true causal effect of citrus intake on the probability of having curly hair.} \\
\\ 
\texttt{Comparing the two states:} \\
\texttt{{-} Without citrus, the probability of curly hair is 47\%.} \\
\texttt{{-} With citrus, the probability of curly hair drops to 15\%.} \\
\\ 
\texttt{Since \$15\textbackslash{}\% < 47\textbackslash{}\%\$, the intervention of consuming citrus results in a lower probability of the outcome (curly hair) compared to not consuming it. Therefore, citrus intake causally decreases the chance of curly hair.} \\
            }
        }
    }
    & & \\ \\

    \theutterance \stepcounter{utterance}  
    & & & \multicolumn{2}{p{0.3\linewidth}}{
        \cellcolor[rgb]{0.9,0.9,0.9}{
            \makecell[{{p{\linewidth}}}]{
                \texttt{\tiny{[GM$|$GM]}}
                \texttt{Yes} \\
            }
        }
    }
    & & \\ \\

    \theutterance \stepcounter{utterance}  
    & & & \multicolumn{2}{p{0.3\linewidth}}{
        \cellcolor[rgb]{0.9,0.9,0.9}{
            \makecell[{{p{\linewidth}}}]{
                \texttt{\tiny{[GM$|$GM]}}
                \texttt{yes} \\
            }
        }
    }
    & & \\ \\

    \theutterance \stepcounter{utterance}  
    & & & \multicolumn{2}{p{0.3\linewidth}}{
        \cellcolor[rgb]{0.9,0.9,0.9}{
            \makecell[{{p{\linewidth}}}]{
                \texttt{\tiny{[GM$|$GM]}}
                \texttt{game\_result = WIN} \\
            }
        }
    }
    & & \\ \\

\end{supertabular}
}

\end{document}
